\thispagestyle{empty}
\begin{center}
\vspace*{3.2cm}
{\Huge\bfseries Physics}\\[6mm]
{\Large A Reasoning-First Companion}\\[4mm]
{\large for JEE and NEET}\\[14mm]
{\color{rulegrey}\rule{0.5\textwidth}{0.9pt}}\\[8mm]
\begin{minipage}{0.76\textwidth}\centering\large\itshape
derive rather than assert: a result you can rebuild beats one you can only recall
\end{minipage}\\[14mm]
{\color{rulegrey}\rule{0.5\textwidth}{0.4pt}}\\[12mm]
{\normalsize
Foundations\ \ The Mathematical Toolkit\\[4mm]
Volume I\ \ Mechanics\\[2mm]
Volume II\ \ Properties of Matter and Thermal Physics\\[2mm]
Volume III\ \ Nuclei and Radioactivity}
\vfill
{\small Aligned to the rationalised NCERT curriculum,\\
JEE Main 2026 and NEET UG 2026.}
\vspace{2cm}
\end{center}
\cleardoublepage

\chapter*{How to Use This Book}
\addcontentsline{toc}{chapter}{How to Use This Book}
\markboth{How to Use This Book}{How to Use This Book}

The book is published as one document in volumes. A short opening chapter, the \emph{Foundations},
collects the mathematics that every volume uses, and its boxes and figures are numbered F.1, F.2 and so
on. Each volume then numbers its own parts and chapters from the beginning, so that a reference within
a volume reads ``Chapter~11'' and one across volumes reads ``Volume~I, Chapter~11''. % stylecheck: ok

Volume~I, \emph{Mechanics}, rests on a single promise: mechanics is one law applied with care. Newton's
second law, $\vF=m\va$, resolved into components and supplemented by two conservation principles,
energy and momentum, accounts for everything from a sliding block to a planet's orbit. Everything else
(the formula for the range of a projectile, for the period of a satellite, for the speed of a rolling
sphere) is bookkeeping built on that foundation. Volume~II, \emph{Properties of Matter and Thermal
Physics}, applies the same mechanics to matter in bulk: how solids and fluids respond to force, how
matter responds to heat, and how both answers follow from molecules in motion. Volume~III, \emph{Nuclei
and Radioactivity}, goes inside the atom: what holds a nucleus together, why some nuclei decay and by
what law, and how the same binding energy powers reactors, the Sun and the clocks that date rocks and
bones. Throughout, the book
derives rather than asserts, because a result you can rebuild beats one you can only recall.

\section*{The conventions of the page}
\begingroup\renewcommand{\thebxcount}{\unskip}

A textbook is read in two ways at once: straight through, and by jumping to the one thing you need
right now. The typography is arranged to serve both. There is no colour code to learn. Instead three
devices mark three kinds of material, and each announces itself by name where it begins.

\begin{defbox}{Statements you are expected to retain}
Material set between two rules is a \emph{definition}, a \emph{result}, a \emph{method} or a
\emph{table}: the load-bearing statements of the subject. A definition fixes the meaning of a symbol; a
result is something that has just been earned by argument and will be used again later; a method is a
procedure to follow; a table collects results for reference. All boxes share one sequence of numbers
within each chapter, so that any of them can be cited precisely.
\end{defbox}

\begin{derivbox}{Arguments you are expected to be able to rebuild}
A hairline rule down the left margin marks a \emph{derivation}, a \emph{worked example}, a
\emph{problem type} or a \emph{challenge}. All four share one property: they are arguments rather than
assertions, and every step is shown, so that you can reproduce the whole thing on blank paper. Worked
examples close with a small open square. A \emph{remark} (an aside on history or an extension beyond
the syllabus) is set in the same way, in smaller type.
\end{derivbox}

\begin{trapbox}{how the two left-hand rules differ}
A lighter hairline, with the opening words in italic as here, carries commentary rather than content:
a caution about a mistake examiners bait, a note on why a result matters beyond the exam hall, a
plain-English gloss on an equation, or a remark on what the syllabus does and does not ask. Skip these
on a first reading and you lose no logic. You will, however, lose most of the reasons.
\end{trapbox}

Illustrations are the one place colour is used, and there it carries information: distinct colours
mark distinct physical quantities, used consistently from one chapter to the next. Every figure is
numbered by chapter and captioned so that it can be understood without the surrounding text.

\endgroup

\section*{On rigour, and on candour}

Numbers in worked examples, problem types and challenge problems have been recomputed independently
rather than copied from a source. Where a popular ``explanation'' is a myth, this book says so and gives the honest
mechanism instead, usually by estimating the size of the effect before accepting the story told about
it. Where a result holds only within an approximation, the approximation is named. Where the full
treatment lies beyond the syllabus, you are told that it exists and told what has replaced it, rather
than being handed a formula that appears to come from nowhere.

That policy costs length and buys trust. A student who has been told once that a textbook was
simplifying has a reason to believe it the next time it says that something is exact.

\section*{How to work through it}

Read a chapter's opening paragraphs before anything else: they state the problem the chapter exists to
solve, and the rest will not cohere without them. Then read the derivations with a pen, reconstructing
each step rather than following it. Attempt every worked example before reading its solution. The
problem types at the end of each chapter collect the recurring shapes that examiners actually set,
with the tell that identifies each one, the trap each one carries, and an exam-style question with a
single correct option; they are a reference to return to, not a list to memorise. Several chapters add
practice problems with answers, and some add challenge problems. The challenges are harder than a
typical examination question, and are meant to be struggled with. The parts of each volume close with a short
retrospective, \emph{Looking Back}, that draws their chapters together; Volume~I ends with one on the
whole of mechanics.

\cleardoublepage
\tableofcontents
\cleardoublepage
